\documentclass[11pt]{article}
\usepackage[T1]{fontenc}
\usepackage[portuguese]{babel}
\usepackage[margin=1in]{geometry}

% Font: PT Sans
\usepackage{paratype}
\renewcommand{\familydefault}{\sfdefault}

\usepackage[most]{tcolorbox}
\usepackage{enumitem}
\usepackage{xcolor}
\usepackage{etoolbox}

% Define Colors
\definecolor{darkblue}{HTML}{00008B}
\definecolor{normalGreen}{HTML}{27AE60}
\definecolor{darkorange}{HTML}{CC8400}
\definecolor{darkred}{HTML}{B22222}
\definecolor{lightgrey}{HTML}{969696}

% Custom Styles for Lists
\setlist[enumerate]{nosep, leftmargin=1.5em, label=\arabic*.}

% --- The Use Case Command ---
\newcommand{\ucspec}[8]{
	\begin{tcolorbox}[
		colframe=lightgrey, 
		colback=white, 
		arc=4pt, 
		boxrule=0.5pt,
		left=12pt, right=12pt, top=12pt, bottom=12pt
		]
		{\Large\bfseries\color{darkblue!60!black} Caso de Uso: #1}
		
		\textcolor{lightgrey!70!black}{\hrule height 0.5pt}
		\vspace{10pt}
		
		\textbf{Descrição:} #2\\[4pt]
		\textbf{Cenários:} #3\\[4pt]
		\textbf{Pré-Condições:} #4\\[4pt]
		\textbf{Pós-Condições:} #5
		
		\vspace{10pt}
		{\bfseries\color{normalGreen} Fluxo Normal:}
		\begin{enumerate}
			#6
		\end{enumerate}
		
		\ifstrempty{#7}{}{
			\vspace{10pt}
			{\bfseries\color{darkorange} Fluxos Alternativos:}
			
			#7
		}
		
		\ifstrempty{#8}{}{
			\vspace{10pt}
			{\bfseries\color{darkred} Fluxos de Exceção:}
			
			#8
		}
	\end{tcolorbox}
}

% Helper for sub-flows
\newcommand{\subflow}[3]{%
	\vspace{6pt}
	\hspace*{10pt} \textbf{Condição:} #1 (Passo #2)
	\begin{enumerate}[label=#2.\arabic*, leftmargin=3.5em]
		#3
	\end{enumerate}
}

\begin{document}
	
	\ucspec
	{Nome do Use Case}
	{Descrição breve do Use Case}
	{Identificar cenários em que o Use Case está presente}
	{Condições necessárias para a execução do Use Case}
	{Condições a satisfazer no fim do Use Case}
	{% Fluxo Normal
		\item 
		\item 
		\item 
		\item 
	}
	{% Fluxos Alternativos
		\subflow{Condição que deu origem ao fluxo alternativo}{1}{
			\item 
			\item 
			\item 
		}
	}
	{% Fluxos de Exceção
		\subflow{Condição que deu origem ao fluxo de exceção}{2}{
			\item 
			\item 
		}
	}

	\ucspec
	{Nome do Use Case}
	{Descrição breve do Use Case}
	{Identificar cenários em que o Use Case está presente}
	{Condições necessárias para a execução do Use Case}
	{Condições a satisfazer no fim do Use Case}
	{% Fluxo Normal
		\item 
		\item 
		\item 
		\item 
	}
	{}
	{}
	
	\ucspec
	{Levantar Dinheiro}
	{Cliente levanta quantia da máquina.}
	{Cenários 1 e 2...}
	{Sistema tem notas.}
	{Cliente tem quantia desejada...}
	{% Fluxo Normal
		\item Cliente apresenta cartão e PIN.
		\item Máquina MB valida acesso e pede operação.
		\item Cliente indica que pretende levantar dada quantia.
		\item Máquina MB pergunta se quer talão.
		\item Cliente responde que não.
		\item Máquina MB devolve cartão, fornece notas e actualiza saldo da conta.
		\item Cliente retira cartão e notas.
	}
	{% Fluxos Alternativos
		\subflow{Cliente quer autenticar-se com reconhecimento facial.}{1}{
			\item Cliente apresenta cartão e pede reconhecimento facial.
			\item Máquina recolhe imagem para validação de acesso.
			\item Regressa a 2.
		}
		\subflow{Cliente quer talão.}{5}{
			\item Cliente responde que sim.
			\item Máquina MB devolve cartão, notas e talão.
			\item Cliente retira cartão, notas e talão e actualiza saldo da conta.
		}
	}
	{% Fluxos de Exceção
		\subflow{O sistema não encontra correspondência para os dados inseridos.}{3}{
			\item Máquina MB avisa sobre PIN inválido e devolve cartão.
			\item Cliente retira cartão.
		}
	}
	
\end{document}